\documentclass[twoside]{article}

\usepackage[preprint]{aistats2027}
\usepackage[round,authoryear]{natbib}
\usepackage{amsmath,amssymb,amsthm}
\usepackage{mathtools}
\usepackage{microtype}
\usepackage{booktabs}
\usepackage{url}

\renewcommand{\bibname}{References}
\renewcommand{\bibsection}{\subsubsection*{\bibname}}
\bibliographystyle{apalike}

\newtheorem{theorem}{Theorem}
\newtheorem{proposition}[theorem]{Proposition}
\newtheorem{corollary}[theorem]{Corollary}
\newtheorem{lemma}[theorem]{Lemma}
\newtheorem{remark}[theorem]{Remark}
\newcommand{\E}{\mathbb E}
\newcommand{\R}{\mathbb R}
\newcommand{\Var}{\operatorname{Var}}
\newcommand{\tr}{\operatorname{tr}}
\newcommand{\rank}{\operatorname{rank}}
\newcommand{\op}{\mathrm{op}}
\newcommand{\cL}{\mathcal L}
\newcommand{\OPT}{\mathrm{OPT}}

\begin{document}
\runningauthor{}

\onecolumn
\aistatstitle{Supplement: Pure Relative Structured Matrix Learning with Square-Root Matvecs}

\aistatsauthor{Pahan Dewasurendra}

\aistatsaddress{Johns Hopkins University}

\appendix

\section{Partial-Trace Covariance}
\label{app:covariance}

This section proves the concentration lemma used in the main paper. We state
a slightly more general accuracy form.

\begin{lemma}[Partial-trace Gaussian covariance]
\label{lem:app-covariance}
Let $C_1,\ldots,C_q\in\R^{n\times m}$ satisfy
\begin{equation}
 K=(\langle C_i,C_j\rangle_F)_{i,j}\preceq I,
 \qquad
 \sum_{i=1}^qC_iC_i^\top\preceq\tau I,
 \label{eq:app-assumptions}
\end{equation}
where $\tau\ge1$. For $g\sim N(0,I_m)$, let
$D(g)=[C_1g,\ldots,C_qg]\in\R^{n\times q}$. There is a universal constant
$C$ such that, for every $\eta\in(0,1)$ and
\begin{equation}
 k\ge \frac{C\tau\log^2(2q)}{\eta^2},
 \label{eq:app-k}
\end{equation}
independent $g_1,\ldots,g_k$ satisfy
\begin{equation}
 \left\|\frac1k\sum_{\ell=1}^k
 D(g_\ell)^\top D(g_\ell)-K\right\|_\op\le\eta
 \label{eq:app-conclusion}
\end{equation}
with probability at least $0.99$.
\end{lemma}

We use the following standard rectangular form of noncommutative Khintchine
\citep{buchholz2001operator}. The Gaussian statement follows from the same
trace-moment proof as the Rademacher statement. It can also be obtained by
self-adjoint dilation.

\begin{lemma}[Rectangular Gaussian Khintchine]
\label{lem:nck}
For fixed real matrices $A_1,\ldots,A_m$ and an integer $p\ge1$,
\begin{align}
 &\left(\E\left\|\sum_{a=1}^m g_aA_a\right\|_{S_{2p}}^{2p}
 \right)^{1/(2p)}\notag\\
 &\quad\le \sqrt{2p-1}\max\left\{
 \left\|\left(\sum_aA_aA_a^\top\right)^{1/2}\right\|_{S_{2p}},\right.\notag\\
 &\hspace{105pt}\left.\left\|\left(\sum_aA_a^\top A_a\right)^{1/2}
 \right\|_{S_{2p}}\right\}.
 \label{eq:nck}
\end{align}
Here $\|\cdot\|_{S_r}$ denotes the Schatten $r$-norm.
\end{lemma}

\begin{proof}[Proof of Lemma~\ref{lem:app-covariance}]
For each input coordinate $a\in[m]$, define
\[
 A_a=[C_1e_a,\ldots,C_qe_a]\in\R^{n\times q}.
\]
Then $D(g)=\sum_ag_aA_a$, and direct calculation gives
\begin{equation}
 \sum_aA_aA_a^\top=\sum_iC_iC_i^\top\preceq\tau I,
 \qquad
 \sum_aA_a^\top A_a=K\preceq I.
 \label{eq:two-variances}
\end{equation}
Both matrices in \eqref{eq:two-variances} have trace $\tr K\le q$.
Consequently,
\begin{align}
 \tr\left(\sum_aA_aA_a^\top\right)^p
 &\le\tau^{p-1}q,\label{eq:left-trace}\\
 \tr K^p&\le q.\label{eq:right-trace}
\end{align}
Apply Lemma~\ref{lem:nck}, use
$\|M\|_\op\le\|M\|_{S_{2p}}$, and take $p\ge\log(2q)$. Since
$q^{1/(2p)}$ is bounded by a universal constant and $\tau\ge1$, we obtain
\begin{equation}
 \left(\E\|D(g)\|_\op^{2p}\right)^{1/(2p)}
 \le C\sqrt{p\tau}.
 \label{eq:D-moment}
\end{equation}

Set $Z=D(g)^\top D(g)$, so $\E Z=K$. We next compute its centered matrix
variance. For a unit vector $\theta\in\R^q$, let
$C_\theta=\sum_i\theta_iC_i$. Since
\[
 \theta^\top Z^2\theta
 =\sum_{i=1}^q\langle C_ig,C_\theta g\rangle^2,
\]
the real Gaussian fourth-moment identity gives
\begin{align}
 \theta^\top\E Z^2\theta
 &=\sum_i\left\{
 \langle C_i,C_\theta\rangle_F^2
 +\tr[(C_i^\top C_\theta)^2]\right\}\notag\\
 &\quad+\sum_i\|C_i^\top C_\theta\|_F^2\notag\\
 &\le\theta^\top K^2\theta
 +2\sum_i\|C_i^\top C_\theta\|_F^2\notag\\
 &=\theta^\top K^2\theta
 +2\tr\left(C_\theta^\top
 \sum_iC_iC_i^\top C_\theta\right)\notag\\
 &\le\theta^\top(K^2+2\tau K)\theta.
 \label{eq:fourth}
\end{align}
The first equality uses
$\E(g^\top Mg)^2=(\tr M)^2+\tr(M^2)+\tr(MM^\top)$.
Because $\E(Z-K)^2=\E Z^2-K^2$, equation \eqref{eq:fourth} implies
\begin{equation}
 \|\E(Z-K)^2\|_\op\le2\tau.
 \label{eq:centered-variance}
\end{equation}

Let $X_\ell=Z_\ell-K$ for independent copies $Z_\ell$. The variance
parameter of $\sum_\ell X_\ell$ is at most $2k\tau$. We also need the
large-deviation parameter
\[
 L=\left(\E\max_{\ell\le k}\|X_\ell\|_\op^2\right)^{1/2}.
\]
For any integer $p\ge\log(2q)$,
\begin{align*}
 \|\|X_\ell\|_\op\|_{L_{2p}}
 &\le \|\|D(g_\ell)\|_\op^2\|_{L_{2p}}+1\\
 &=\|\|D(g_\ell)\|_\op\|_{L_{4p}}^2+1
 \le Cp\tau.
\end{align*}
The elementary maximum bound gives
\begin{align*}
 L
 &\le\left(\E\max_{\ell\le k}\|X_\ell\|_\op^{2p}
 \right)^{1/(2p)}\\
 &\le k^{1/(2p)}\max_{\ell\le k}
 \|\|X_\ell\|_\op\|_{L_{2p}}.
\end{align*}
Taking $p=\lceil\log(2qk)\rceil$ yields
\begin{equation}
 L\le C\tau\log(2qk).
 \label{eq:large-deviation}
\end{equation}

The expected-norm theorem for independent centered random matrices
\citep[Theorem I]{tropp2015expected}, applied in dimension $q$, now gives
\begin{multline}
 \left(\E\left\|\frac1k\sum_{\ell=1}^kX_\ell\right\|_\op^2
 \right)^{1/2}\\
 \le C\sqrt{\frac{\tau\log(2q)}{k}}
 +\frac{C\tau\log(2q)\log(2qk)}{k}.
 \label{eq:tropp-application}
\end{multline}

Let $k_0=C_0\tau\log^2(2q)/\eta^2$. Since $\tau$ may be replaced by
$\min\{\tau,q\}$, we may assume $\tau\le q$. At $k=k_0$,
\[
 \log(2qk)=O(\log(2q)+\log(1/\eta)).
\]
The first term of \eqref{eq:tropp-application} is at most $C\eta/\sqrt{\log(2q)}$.
For the second term, use
$\eta\log(1/\eta)\le1/e$. Increasing $C_0$ makes the root second moment at
most $\eta/10$. For $k\ge k_0$, the first term decreases and
$\log(2qk)/k$ also decreases. Markov's inequality proves
\eqref{eq:app-conclusion} with probability at least $0.99$.
\end{proof}

\section{Upper Bound and Amplification}
\label{app:upper}

We first record basis invariance and observability.

\begin{lemma}[Basis invariance]
\label{lem:basis}
The partial traces
$S_L=\sum_iB_iB_i^\top$ and $S_R=\sum_iB_i^\top B_i$ do not depend on the
choice of Frobenius-orthonormal basis of $\cL$.
\end{lemma}

\begin{proof}
Any two orthonormal bases are related by an orthogonal matrix $U$. If
$B'_i=\sum_jU_{ij}B_j$, then
\[
 \sum_iB'_iB_i'^\top
 =\sum_{j,k}(U^\top U)_{jk}B_jB_k^\top
 =\sum_jB_jB_j^\top.
\]
The right partial trace is identical after transposition.
\end{proof}

\begin{lemma}[Observable objective]
\label{lem:observable}
Let $P=VV^\top$, where $V\in\R^{n\times s}$ has orthonormal columns. The
responses $A^\top V$ and $AG$ determine
\[
 \|P(A-B)\|_F^2+k^{-1}\|Q(A-B)G\|_F^2
\]
for every known $B$.
\end{lemma}

\begin{proof}
The left responses determine
$PA=V(A^\top V)^\top$. Therefore $PAG$ and
$QAG=AG-PAG$ are known. All terms involving $B$ are known from the basis.
\end{proof}

\begin{lemma}[Offset variance]
\label{lem:app-offset}
Suppose $\sum_iC_iC_i^\top\preceq\lambda I$. For fixed $M$, define
\[
 \xi_i=\frac1k\sum_{\ell=1}^k
 \{\langle Mg_\ell,C_ig_\ell\rangle-\langle M,C_i\rangle_F\}.
\]
Then $\E\|\xi\|_2^2\le2\lambda\|M\|_F^2/k$.
\end{lemma}

\begin{proof}
For $N=M^\top C_i$, the Gaussian identity
$\Var(g^\top Ng)=\tr(N^2)+\tr(NN^\top)$ gives
\[
 \Var(\langle Mg,C_ig\rangle)\le2\|M^\top C_i\|_F^2.
\]
Summing over $i$ and using independence across $\ell$ gives
\begin{align*}
 \E\|\xi\|_2^2
 &\le\frac2k\sum_i\|M^\top C_i\|_F^2\\
 &=\frac2k\tr\left(M^\top\sum_iC_iC_i^\top M\right)
 \le\frac{2\lambda}{k}\|M\|_F^2.
\end{align*}
\end{proof}

\begin{theorem}[One-orientation upper bound]
\label{thm:app-upper}
Fix $s$ and let $\lambda=\lambda_{s+1}^L>0$. For every
$\gamma\in(0,1)$, if
\begin{equation}
 k\ge C\left[
 \max\{1,\lambda\}\log^2(2q)+\frac{\lambda}{\gamma}
 \right],
 \label{eq:app-upper-k}
\end{equation}
then the estimator in the main paper satisfies
\begin{equation}
 \|\widehat B-B_\star\|_F^2\le\gamma\OPT^2
 \label{eq:parameter-target}
\end{equation}
with probability at least $19/20$.
\end{theorem}

\begin{proof}
Let $C_i=QB_i$ and $D(g)=[C_1g,\ldots,C_qg]$. Its Gram matrix
$K=(\langle C_i,C_j\rangle_F)$ satisfies $K\preceq I$, while
\[
 \sum_iC_iC_i^\top=QS_LQ\preceq\lambda I.
\]
Set $\tau=\max\{1,\lambda\}$. The Hessian of the objective in the
orthonormal coordinates of $\cL$ is
\begin{align*}
 H&=H_P+\frac1k\sum_{\ell=1}^kD(g_\ell)^\top D(g_\ell),\\
 (H_P)_{ij}&=\langle PB_i,PB_j\rangle_F.
\end{align*}
Because $H_P+K=I$, Lemma~\ref{lem:app-covariance} with $\eta=1/2$ implies
\begin{equation}
 H\succeq\frac12I
 \label{eq:hessian-event}
\end{equation}
except with probability $1/100$.

Write $R=A-B_\star$. At $B_\star$, the $i$th normal-equation offset is
\begin{align*}
 h_i
 &=\langle PR,PB_i\rangle_F
 +\frac1k\sum_{\ell=1}^k
 \langle QRg_\ell,QB_ig_\ell\rangle.
\end{align*}
Orthogonality gives
\[
 \langle PR,PB_i\rangle_F+\langle QR,QB_i\rangle_F
 =\langle R,B_i\rangle_F=0.
\]
Thus $h$ is the centered vector in Lemma~\ref{lem:app-offset} with $M=QR$.
In particular,
\begin{equation}
 \E\|h\|_2^2\le\frac{2\lambda}{k}\OPT^2.
 \label{eq:h-bound}
\end{equation}
By Markov's inequality and a sufficiently large constant in
\eqref{eq:app-upper-k},
\begin{equation}
 \|h\|_2^2\le\frac\gamma4\OPT^2
 \label{eq:offset-event}
\end{equation}
except with probability $1/25$.

Let $d$ be the coefficient vector of $\widehat B-B_\star$. On
\eqref{eq:hessian-event}, the objective is strictly convex and its normal
equation is $Hd=h$. Therefore
\[
 \|\widehat B-B_\star\|_F^2=\|d\|_2^2
 \le4\|h\|_2^2\le\gamma\OPT^2.
\]
A union bound gives success probability at least
$1-1/100-1/25=19/20$.
\end{proof}

If $\lambda=0$, then $QB_i=0$ for all $i$. The exact part of the objective
has Hessian $I$ and recovers $B_\star$ without random queries. The profile in
the main paper follows by taking $\gamma=2\epsilon$ in
Theorem~\ref{thm:app-upper} and applying
\[
 \|A-\widehat B\|_F^2
 =\OPT^2+\|\widehat B-B_\star\|_F^2.
\]

We next prove amplification without estimating $\OPT$ or any residual norm.

\begin{lemma}[Metric median amplification]
\label{lem:median}
Let $\widehat B_1,\ldots,\widehat B_r$ be independent candidates, where $r$
is odd and each candidate satisfies
\begin{equation}
 \|\widehat B_i-B_\star\|_F^2\le\gamma\OPT^2
 \label{eq:good}
\end{equation}
with probability at least $19/20$. Let $h=(r+1)/2$. For each $i$, let
$\rho_i$ be the $h$th smallest value among
$\{\|\widehat B_i-\widehat B_j\|_F:j\in[r]\}$. Return a candidate minimizing
$\rho_i$. There is a universal $C$ such that $r\ge C\log(1/\delta)$ implies
\begin{equation}
 \|\widehat B-B_\star\|_F^2\le9\gamma\OPT^2
 \label{eq:median-conclusion}
\end{equation}
with probability at least $1-\delta$.
\end{lemma}

\begin{proof}
A Chernoff bound shows that at least $h$ candidates satisfy \eqref{eq:good}
with probability at least $1-\delta$. Condition on this event. Every good
candidate has at least $h$ good candidates within distance
$2\sqrt\gamma\OPT$, so its radius is at most this value. The selected
candidate has radius no larger. Its radius ball contains at least $h$
candidates, and the good set also has at least $h$ members. Since $2h>r$,
the two sets intersect. If $\widehat B_j$ is in their intersection, then
\[
 \|\widehat B-B_\star\|_F
 \le\|\widehat B-\widehat B_j\|_F
 +\|\widehat B_j-B_\star\|_F
 \le3\sqrt\gamma\OPT.
\]
\end{proof}

Take $\gamma=2\epsilon/9$ in Lemma~\ref{lem:median}. Then
\[
 \|A-\widehat B\|_F^2
 \le(1+2\epsilon)\OPT^2
 \le(1+\epsilon)^2\OPT^2.
\]
The same exact responses $A^\top V$ are reused for all candidates. Only the
Gaussian stage is repeated, proving the confidence statement in the main
paper.

\section{Lower Bounds}
\label{app:lower}

\subsection{Exact-recovery dimension bound}

\begin{proposition}
\label{prop:dimension-lower}
For $q=r^2$, some $q$-dimensional linear family requires at least $r$ adaptive
two-sided queries for any finite multiplicative approximation guarantee with
positive worst-case success probability.
\end{proposition}

\begin{proof}
Take $\cL=\R^{r\times r}$ and draw $A$ from a nondegenerate Gaussian
distribution on this space. Fix a deterministic adaptive algorithm using
$t<r$ queries. Conditional on any transcript up to a given round, the next
query vector is fixed. Its response gives at most $r$ scalar linear
constraints on the $r^2$ entries of $A$, whether the orientation is $A$ or
$A^\top$. Sequential Gaussian conditioning shows that after $t$ rounds the
conditional law remains Gaussian on an affine space of dimension at least
$r^2-tr>0$. It is non-atomic. No transcript-measurable output equals $A$ with
positive probability.

For every $A\in\cL$, the optimum residual is zero. Any finite multiplicative
guarantee must output $A$ exactly. Therefore every deterministic algorithm
has zero average success under the Gaussian prior. Yao's principle rules out
a randomized algorithm with positive worst-case success.
\end{proof}

\subsection{One-parameter accuracy bound}

We give the deferred-decisions argument and constants for the one-parameter
proposition in the main paper. Write
$G\sim\operatorname{GOE}_d$ when the diagonal entries of $G$ are independent
$N(0,2)$ variables and its entries strictly above the diagonal are
independent $N(0,1)$ variables, with symmetry below the diagonal.

\begin{lemma}[Adaptive GOE residual]
\label{lem:goe-residual}
Fix a deterministic adaptive algorithm that queries a symmetric matrix by
matvecs. After $t$ linearly independent queries to
$G\sim\operatorname{GOE}_d$, there is a transcript-measurable orthogonal
matrix $U_t$ such that, conditionally on the transcript,
\begin{equation}
 U_t^\top G U_t=\Delta_t+
 \begin{bmatrix}
  0_{t\times t}&0\\
  0&H_t
 \end{bmatrix},
 \qquad H_t\sim\operatorname{GOE}_{d-t},
 \label{eq:goe-residual}
\end{equation}
where $\Delta_t$ is transcript measurable and $H_t$ is independent of the
transcript.
\end{lemma}

\begin{proof}
Repeated queries and components in the span of previous queries reveal no new
information, so Gram--Schmidt reduces to orthonormal query directions. We
prove the claim by induction. It is immediate for $t=0$. Suppose it holds
after $t$ queries. In the coordinates of \eqref{eq:goe-residual}, the next
new query has a nonzero component in the last $d-t$ coordinates. Conditional
on the transcript, rotate only those coordinates so that its normalized new
component is the first coordinate. Rotational invariance leaves $H_t$ GOE.
The response reveals the first row and column of this rotated $H_t$. By the
independent-entry definition of GOE, its remaining principal block is an
independent $\operatorname{GOE}_{d-t-1}$ matrix. Absorb every revealed entry
into $\Delta_{t+1}$. This proves the induction, including when each query is
chosen from all earlier responses.
\end{proof}

\begin{proposition}[One-parameter accuracy lower bound]
\label{prop:app-one-parameter}
For all sufficiently small $\epsilon>0$, some one-dimensional linear family
requires $\Omega(1/\sqrt\epsilon)$ adaptive two-sided queries for success
probability at least $2/3$ at relative error $1+\epsilon$.
\end{proposition}

\begin{proof}
Set $c=1/100$, let $d=\lfloor c/\sqrt\epsilon\rfloor$, and take
\[
 \cL=\operatorname{span}\{I_d/\sqrt d\},
 \qquad A\sim\operatorname{GOE}_d.
\]
The input is symmetric, so a transpose query reveals the same response as a
forward query. It is enough by Yao's principle to fix a deterministic
adaptive algorithm. Its output is $\widehat\beta I_d/\sqrt d$, while the
optimal coefficient is $\beta_\star=\tr(A)/\sqrt d$.

After $t\le d/2$ informative queries, Lemma~\ref{lem:goe-residual} shows that
the unknown part of $\beta_\star$, conditionally on the transcript, is
\[
 \frac{\tr(H_t)}{\sqrt d}\sim
 N\left(0,\frac{2(d-t)}d\right).
\]
Its conditional variance is at least one. The density of a Gaussian with
variance at least one is everywhere at most $1/\sqrt{2\pi}$. Hence, for every
transcript-measurable estimate $\widehat\beta$ and every $a>0$,
\begin{equation}
 \Pr\bigl\{|\widehat\beta-\beta_\star|\le a
       \mid\text{transcript}\bigr\}
 \le\frac{2a}{\sqrt{2\pi}}.
 \label{eq:goe-small-ball}
\end{equation}

We have $\E\|A\|_F^2=d(d+1)$, so Markov's inequality gives the event
\begin{equation}
 \OPT^2\le\|A\|_F^2\le10d(d+1)
 \label{eq:goe-norm-event}
\end{equation}
with probability at least $0.9$. For $\epsilon\le1$, a successful
$(1+\epsilon)$ approximation must obey
\begin{align*}
 (\widehat\beta-\beta_\star)^2
 &\le\bigl((1+\epsilon)^2-1\bigr)\OPT^2\\
 &\le3\epsilon\OPT^2.
\end{align*}
On \eqref{eq:goe-norm-event}, the right side is at most
$30\epsilon d(d+1)\le60c^2$. Apply \eqref{eq:goe-small-ball} with
$a=\sqrt{60}c$. The total success probability is at most
\[
 \frac1{10}+\frac{2\sqrt{60}c}{\sqrt{2\pi}}<\frac16.
\]
For sufficiently small $\epsilon$, we also have
$d\ge c/(2\sqrt\epsilon)$. Therefore success probability $2/3$ requires
more than $d/2=\Omega(1/\sqrt\epsilon)$ queries under this input
distribution. Yao's principle gives the claimed worst-case lower bound for
randomized algorithms.
\end{proof}

\subsection{Joint accuracy bound}

We state the source theorem in the form needed here. It is Theorem 2 of
\citet{amsel2026fixed}, together with the dimension choice in its proof.

\begin{proposition}[Fixed-sparsity Wishart lower bound]
\label{prop:source-lower}
Fix $\alpha\in(0,1)$. There are constants $a,b,c_0,c_1>0$ such that, for
every integer $u\ge1$ and $\epsilon<c_0$, one can choose
\begin{equation}
 a\frac{u}{\epsilon}\le d\le b\frac{u}{\epsilon}+1
 \label{eq:d-range}
\end{equation}
and a symmetric Wishart distribution on $d\times d$ matrices with this
property. For every mask having between $\alpha u$ and $u$ entries in every
row and column, any possibly randomized adaptive algorithm using fewer than
$c_1u/\epsilon$ matvec queries has probability at most $1/25$ of returning a
$(1+\epsilon)$ fixed-mask approximation.
\end{proposition}

The source theorem is stated for forward queries. Its hard matrix is
symmetric, so allowing transpose queries does not enlarge the transcript.

\begin{theorem}[Joint lower bound for linear families]
\label{thm:joint-lower}
There are universal constants $c,c_0,c_2>0$ such that, whenever
$\epsilon<c_0$ and $q\epsilon\ge c_2$, some $q$-dimensional linear family
requires at least $c\sqrt{q/\epsilon}$ possibly adaptive two-sided queries
for constant-probability $(1+\epsilon)$ approximation.
\end{theorem}

\begin{proof}
For each $u$ and $d$ in Proposition~\ref{prop:source-lower}, choose a cyclic
$u$-regular mask $S\in\{0,1\}^{d\times d}$. This exists after reducing
$c_0$ if necessary so that $d\ge u$. The coordinate family
\[
 \cL_0=\{B:B=S\circ B\}
\]
has dimension
\begin{equation}
 q_0=du=\Theta(u^2/\epsilon).
 \label{eq:q0}
\end{equation}
Proposition~\ref{prop:source-lower} gives query lower bound
\begin{equation}
 \Omega(u/\epsilon)=\Omega(\sqrt{q_0/\epsilon}).
 \label{eq:source-rate}
\end{equation}

Given $q$ with $q\epsilon$ above a large enough constant, choose
$u=\lfloor c_3\sqrt{q\epsilon}\rfloor$ for a sufficiently small universal
$c_3$. Equations \eqref{eq:d-range} and \eqref{eq:q0} ensure
$c_4q\le q_0\le q$ for a universal $c_4>0$. Rounding changes only constants.

Let $p=q-q_0$. On a disjoint $p\times p$ block, let $\mathcal D_p$ be the
diagonal coordinate family, and define
\[
 \cL=\{B_0\oplus D:B_0\in\cL_0,\ D\in\mathcal D_p\}.
\]
This family has dimension exactly $q$. Embed a hard input as
$A'=A_0\oplus0$. Its best approximation in $\cL$ is
$(S\circ A_0)\oplus0$, and its optimum residual equals the original one.

Suppose an algorithm succeeds for $\cL$. A query to $A'$ with vector
$(x,z)$ returns $(A_0x,0)$, and a transpose query is identical. Thus every
query can be simulated using one query to $A_0$. The algorithm's output has
the form $\widehat B_0\oplus\widehat D$. Success implies
\begin{align*}
 \|A_0-\widehat B_0\|_F^2
 &\le\|A_0-\widehat B_0\|_F^2+\|\widehat D\|_F^2\\
 &\le(1+\epsilon)^2
 \|A_0-S\circ A_0\|_F^2.
\end{align*}
Hence restriction to the first block solves the source problem. By
\eqref{eq:source-rate} and $q_0\ge c_4q$, the required number of queries is
$\Omega(\sqrt{q/\epsilon})$.
\end{proof}

\section{Profiles for Basic Families}
\label{app:examples}

\paragraph{Coordinate subspaces.}
For a binary mask $S$, use basis matrices $E_{ij}$ at its nonzero locations.
Then
\begin{align*}
 S_L&=\operatorname{diag}(d_1^{\rm row},\ldots,d_n^{\rm row}),\\
 S_R&=\operatorname{diag}(d_1^{\rm col},\ldots,d_m^{\rm col}).
\end{align*}
The spectral split exactly queries selected high-degree rows or columns and
uses a Gaussian fit on the remainder.

\paragraph{Full rectangular block.}
For all matrices supported on an $a\times b$ block,
$S_L=bI_a$ on the row block and $S_R=aI_b$ on the column block. Taking
$s=a$ in the left profile or $s=b$ in the right profile recovers the block
exactly in $\min\{a,b\}$ queries. For $a=b=\sqrt q$, this matches
Proposition~\ref{prop:dimension-lower}.

\paragraph{Row star.}
For a single active row with $q$ independent entries, $S_L=qee^\top$. One
left query in direction $e$ observes the entire structured component and
gives exact recovery.

\paragraph{Diagonal family.}
For the $q\times q$ diagonal family, $S_L=S_R=I_q$. Taking $s=0$ gives
$O(\log^2(2q)+1/\epsilon)$ queries. The specialized rowwise inverse-Wishart
analysis removes the logarithmic term \citep{amsel2026fixed}. This example
shows that the universal profile is adaptive to diffuse leverage, even though
the general matrix-concentration proof is not logarithmically sharp.

\paragraph{Circulant family.}
Let $P_j$ be the $j$th cyclic shift permutation. The matrices
$B_j=P_j/\sqrt n$ form a Frobenius-orthonormal basis, and
\[
 \sum_{j=0}^{n-1}B_jB_j^\top
 =\sum_{j=0}^{n-1}\frac1n I=I.
\]
Thus $s=0$ gives $O(\log^2(2n)+1/\epsilon)$ queries.

\paragraph{Toeplitz family.}
For offset $j\in\{-(n-1),\ldots,n-1\}$, let $T_j$ be the binary matrix on
the $j$th diagonal and set $B_j=T_j/\sqrt{n-|j|}$. These matrices form an
orthonormal basis. Their left partial trace is diagonal. At row $i$ its entry
is
\begin{align*}
 [S_L]_{ii}
 &=\sum_{j=0}^{n-i}\frac1{n-j}
   +\sum_{h=1}^{i-1}\frac1{n-h}\\
 &=H_n-H_{i-1}+H_{n-1}-H_{n-i}.
\end{align*}
The maximum is $H_n$, attained at an endpoint. The right partial trace is the
reflection of the same diagonal. Taking $s=0$ therefore gives
$O(H_n[\log^2(4n)+1/\epsilon])$ queries.

\paragraph{Known-graph symmetric precision family.}
For a graph $G=(V,E)$ on $n$ vertices, take basis $E_{ii}$ for each vertex
and $(E_{ij}+E_{ji})/\sqrt2$ for each edge. This is Frobenius orthonormal and
\[
 S_L=S_R=\operatorname{diag}
 \left(1+\frac{\deg(1)}2,\ldots,1+\frac{\deg(n)}2\right).
\]
The no-split profile is controlled by $1+\Delta/2$. More generally, querying
the $s$ highest-degree coordinate directions exactly leaves residual leverage
$1+d_{s+1}/2$, where $d_1\ge\cdots\ge d_n$ are the degrees.

\section{Protocol for the Profile Illustration}
\label{app:experiment}

Figure 1 in the main paper uses $32\times32$ circulant, Toeplitz, and path
precision families with their normalized bases above. For each family, NumPy
generator seed 20260729 produces a standard Gaussian matrix. We subtract its
exact Frobenius projection onto the family and normalize the remaining
residual $R$ to $\|R\|_F=1$. The structured optimum is set to zero because the
coefficient error induced by an orthogonal residual is translation invariant.

For each $k\in\{2,3,4,6,8,12,16,24,32\}$, we draw 120 independent standard
Gaussian matrices $G\in\mathbb R^{32\times k}$ and solve
\[
 \widehat B\in\arg\min_{B\in\cL}\|(R-B)G\|_F^2
\]
with NumPy's dense least-squares solver. The reported statistic is
$\|\widehat B\|_F^2/\|R\|_F^2$, which equals squared excess over the optimum
by orthogonality. Solid curves are medians and dotted curves are 90th
percentiles. The supplementary file \texttt{reproduce\_profile.py} generates
the figure using NumPy 2.2.6 and Matplotlib 3.10.9. It uses one CPU process,
no external data, and completes in under one minute on a conventional compute
node.

\bibliography{references}

\end{document}
